\documentclass[conference]{IEEEtran}
\IEEEoverridecommandlockouts

\usepackage{cite}
\usepackage{amsmath,amssymb,amsfonts}
\usepackage{graphicx}
\usepackage{textcomp}
\usepackage{xcolor}
\usepackage{booktabs}
\usepackage{array}
\usepackage{url}

\def\BibTeX{{\rm B\kern-.05em{\sc i\kern-.025em b}\kern-.08em
    T\kern-.1667em\lower.7ex\hbox{E}\kern-.125emX}}

\begin{document}

\title{Waggle: Graph-Backed Persistent Memory for AI Agents\\
with Recursive Memory Context Assembly}

\author{
\IEEEauthorblockN{Abhigyan Shekhar}
\IEEEauthorblockA{Independent Research\\
\texttt{waggle-mcp} open-source project}
}

\maketitle

\begin{abstract}
AI agents lose all conversational context when the context window closes.
Existing approaches---raw context dumps, flat retrieval-augmented generation (RAG),
and note storage---fail to surface relational structure, contradictions, or superseded
decisions at scale.
We present \textbf{Waggle}, a local-first persistent memory system built on a typed
knowledge graph, and \textbf{Recursive Memory Context Assembly (RMCA)}, an algorithm
that assembles compact, high-signal context packs from that graph via query decomposition,
multi-lane retrieval, typed-edge expansion, conflict resolution, and token-budget compression.

We evaluate across three complementary regimes.
\textit{(1) Session retrieval:} On LongMemEval-S (500 questions), \texttt{graph\_raw}
achieves \textbf{89.0\% Exact@5} and \textbf{97.4\% R@5} with \texttt{all-MiniLM-L6-v2};
a held-out split confirms no overfitting (+1.1\,pp gap).
\textit{(2) Synthetic context assembly:} On five RLM-style memory task families,
RMCA achieves perfect scores on pairwise conflict reasoning and ContextReset
(session-boundary state reconstruction) while flat baselines score 0.0; linear aggregation
remains hard for all methods.
\textit{(3) LLM answer quality:} Groq \texttt{llama-3.3-70b-versatile} achieves
EM=1.0, F1=1.0, and hallucination rate=0.0 when given RMCA context packs on pairwise
conflict tasks, consistent across 3 seeds and 2 scales; all flat baselines fail.

Ablation studies identify query decomposition as the only currently causally isolated
load-bearing component; graph expansion and conflict resolution do not yet show a
measurable delta due to benchmark construction limitations.
Synthetic benchmark numbers are not numerically comparable to the RLM paper, which
evaluates real long-context datasets with frontier models.
We release code, benchmark generators, and result artifacts with the paper.
\end{abstract}

\begin{IEEEkeywords}
agent memory, knowledge graph, retrieval-augmented generation, context assembly,
persistent memory, LLM agents, conflict resolution
\end{IEEEkeywords}

\section{Introduction}

AI coding assistants and conversational agents are stateless by design: each session begins
with an empty context window. Practitioners work around this by pasting prior conversation
summaries, maintaining hand-curated notes, or relying on the model to infer context from
the current session alone. None of these approaches scale.

The core problem has three dimensions:
\begin{enumerate}
\item \textbf{Volume.} A long-running project accumulates thousands of decisions, constraints,
and implementation facts. No context window can hold them all.
\item \textbf{Structure.} Decisions depend on constraints; constraints contradict earlier
preferences; implementations derive from decisions. Flat text storage discards this relational
structure.
\item \textbf{Staleness.} Decisions get reversed. Constraints get relaxed. A memory system
that cannot distinguish current from superseded information will mislead the agent.
\end{enumerate}

Retrieval-augmented generation (RAG)~\cite{lewis2020rag} addresses volume but not structure
or staleness. GraphRAG-style systems~\cite{edge2024graphrag} add community-detection graphs
over document chunks but operate on pre-computed summaries and do not model temporal validity
or contradiction. Episodic memory systems~\cite{packer2023memgpt,park2023generative} store
conversation turns but do not extract typed relational structure.

We make the following contributions:
\begin{itemize}
\item \textbf{Waggle}, a local-first persistent memory system built on a typed knowledge graph
with temporal validity, contradiction handling, and a git-inspired snapshot vocabulary.
\item \textbf{RMCA (Recursive Memory Context Assembly)}, an algorithm that assembles compact
context packs from the graph via query decomposition, multi-lane retrieval, typed-edge
expansion, conflict resolution, and token-budget compression.
\item \textbf{Empirical evaluation} on LongMemEval-S (500 questions) and five synthetic task
families, with held-out splits, LLM-backed answer-level evaluation, and explicit ablation
studies.
\item \textbf{Open-source implementation} as an MCP-compatible server (\texttt{waggle-mcp})
deployable with a single \texttt{pipx install}.
\end{itemize}

\section{Background and Related Work}

\subsection{Context Window Limitations}

Transformer-based LLMs have fixed context windows. While window sizes have grown
(4K $\to$ 128K $\to$ 1M tokens), the fundamental problem remains: information outside the
window is inaccessible, and filling the window with raw history is expensive and degrades
attention quality at long ranges~\cite{hsieh2024ruler}.

\subsection{Retrieval-Augmented Generation}

RAG~\cite{lewis2020rag} retrieves relevant text chunks by semantic similarity and injects
them into the prompt. It addresses volume but has three structural weaknesses relevant to
agent memory: flat ranking cannot distinguish current from superseded decisions; chunks are
independent so dependency structure is lost; and two contradictory chunks are returned with
equal weight.

\subsection{GraphRAG and Knowledge Graph Retrieval}

Microsoft GraphRAG~\cite{edge2024graphrag} builds a community-detection graph over document
chunks and retrieves community summaries. This improves coverage for broad queries but
pre-computes summaries offline and does not model temporal validity or agent-specific
relational types (decisions, constraints, contradictions).

\subsection{Episodic and Working Memory for Agents}

MemGPT~\cite{packer2023memgpt} introduces a tiered memory architecture with explicit paging
between in-context and external storage. Generative Agents~\cite{park2023generative} use a
memory stream with recency, importance, and relevance scoring. These systems focus on
retrieval of episodic facts; they do not model typed relational edges or contradiction
resolution as first-class operations.

\subsection{Recursive Language Models}

The RLM framework~\cite{zhang2026rlm} proposes externalising long context into an environment
and interacting with it through decomposition and targeted retrieval. Waggle's RMCA algorithm
is directly inspired by this decomposition-retrieval loop, applied to a persistent typed graph
rather than a document corpus. The RLM paper reports that base models score $<$0.1\% F1 on
OOLONG-Pairs while structured recursive assembly recovers to 23--58\% F1, a qualitative
pattern consistent with our findings.

\section{System Design}

\subsection{Memory Graph}

Waggle stores memory in a persistent typed knowledge graph $G = (V, E)$.

\textbf{Node types:} \texttt{fact}, \texttt{entity}, \texttt{concept}, \texttt{preference},
\texttt{decision}, \texttt{question}, \texttt{note}.

\textbf{Edge types:} \texttt{relates\_to}, \texttt{contradicts}, \texttt{depends\_on},
\texttt{part\_of}, \texttt{updates}, \texttt{derived\_from}, \texttt{similar\_to}.

Every node carries temporal validity fields \texttt{valid\_from} and \texttt{valid\_to}.
Nodes whose \texttt{valid\_to} has passed are excluded from default queries. The
\texttt{resolve\_conflict} operation accepts a \texttt{winner} node ID and sets the losing
node's \texttt{valid\_to} to the current time, making supersession explicit and reversible.

The graph is backed by SQLite with WAL mode for local-first deployments, with an optional
Neo4j backend for remote or multi-client deployments.

\subsection{Embeddings}

Waggle uses \texttt{sentence-transformers}~\cite{reimers2019sbert} for local embedding with
no external API dependency. The default model is \texttt{all-MiniLM-L6-v2} ($\approx$420\,MB,
384-dimensional embeddings). A deterministic SHA-256 fallback is available for offline or
resource-constrained environments.

\subsection{Retrieval Modes}

Three retrieval lanes are available:
\begin{itemize}
\item \textbf{\texttt{graph}}: semantic similarity over node embeddings, with typed-edge expansion.
\item \textbf{\texttt{hybrid}}: graph retrieval fused with BM25 lexical scoring.
\item \textbf{\texttt{verbatim}}: semantic search over the raw conversation transcript store.
\end{itemize}

\subsection{MCP Tool Surface}

Waggle exposes its functionality as an MCP (Model Context Protocol) server.
Table~\ref{tab:tools} lists the six core tools used in normal agent operation.

\begin{table}[htbp]
\caption{Waggle MCP Tool Surface}
\label{tab:tools}
\centering
\begin{tabular}{ll}
\toprule
\textbf{Tool} & \textbf{Purpose} \\
\midrule
\texttt{observe\_conversation} & Ingest a turn; extract nodes and edges \\
\texttt{query\_graph} & Retrieve a subgraph for a query \\
\texttt{prime\_context} & Hydrate context at session start \\
\texttt{build\_context} & Run RMCA for a full context pack \\
\texttt{graph\_diff} & Show what changed recently \\
\texttt{store\_node} / \texttt{store\_edge} & Structured atomic writes \\
\bottomrule
\end{tabular}
\end{table}

\section{Recursive Memory Context Assembly (RMCA)}

RMCA is the algorithm behind \texttt{build\_context}. It takes a query $q$, the memory graph
$G$, the transcript store $T$, a token budget $B$, expansion depth $d$, and maximum subquery
count $k$, and returns a structured context pack $C$ with
$\widehat{\text{tok}}(C) \leq B \times 1.15$.

\subsection{Algorithm}

\textbf{Step 1 --- Decompose.} The query is classified as project/coding or generic.
A set of up to $k$ targeted subqueries is generated, each with a priority weight and a set
of retrieval modes. For project queries, subqueries target: recent decisions, unfinished
tasks, constraints and rejected directions, implementation details, conflicts, and the
original query. For generic queries, subqueries target: the original query, recent relevant
facts, related decisions, contradictions, and transcript evidence.

\textbf{Step 2 --- Retrieve.} Each subquery is issued against the graph and transcript store
using its assigned retrieval modes. Results are accumulated into a working hit set $H$.

\textbf{Step 3 --- Expand.} The top-5 nodes in $H$ by score are used as seeds. For each of
the top-3 seeds, 1--2 hop graph expansion is performed along typed edges (\texttt{updates},
\texttt{contradicts}, \texttt{depends\_on}, \texttt{derived\_from}, \texttt{part\_of}).
Expansion along \texttt{relates\_to} and \texttt{similar\_to} is excluded to avoid token
budget inflation.

\textbf{Step 4 --- Resolve.} All edges in $H$ are inspected. Nodes connected by
\texttt{updates} edges have their scores penalised ($\times 0.3$) and the updating node is
boosted ($+0.15$). Nodes connected by \texttt{contradicts} edges are recorded as conflict
pairs. Nodes with expired \texttt{valid\_to} are penalised ($\times 0.2$).

\textbf{Step 5 --- Deduplicate.} For each unique \texttt{node\_id}, only the highest-scored
copy is retained.

\textbf{Step 6 --- Rank.} Nodes are sorted by: superseded status (non-superseded first),
score (descending), label (ascending for ties).

\textbf{Step 7 --- Compress.} Nodes are packed into the context string $C$ in priority order:
decisions, constraints, implementation, unfinished, conflicts, superseded, evidence. Packing
stops when the token budget $B \times 1.15$ is reached.

\textbf{Step 8 --- Format.} A structured header is prepended. Provenance and conflict
annotations are attached.

\subsection{Comparison with Top-$k$ RAG}

Table~\ref{tab:rag_compare} summarises the key differences between RMCA and standard
top-$k$ RAG.

\begin{table}[htbp]
\caption{RMCA vs.\ Top-$k$ RAG}
\label{tab:rag_compare}
\centering
\begin{tabular}{p{2.0cm}p{2.5cm}p{2.5cm}}
\toprule
\textbf{Dimension} & \textbf{Top-$k$ RAG} & \textbf{RMCA} \\
\midrule
Query strategy & Single embedding query & Decomposed into $\leq k$ targeted subqueries \\
Index structure & Flat chunk index & Typed knowledge graph \\
Graph expansion & None & 1--2 hop traversal along typed edges \\
Conflict handling & None & Explicit \texttt{contradicts}/\texttt{updates} detection \\
Staleness handling & None & Temporal validity (\texttt{valid\_to}) with score penalisation \\
Token budget & Fixed top-$k$ & Configurable budget with priority-ordered packing \\
Evidence provenance & Chunk source & Verbatim transcript lane + node provenance \\
\bottomrule
\end{tabular}
\end{table}

\subsection{Comparison with GraphRAG}

\begin{table}[htbp]
\caption{RMCA vs.\ GraphRAG}
\label{tab:graphrag_compare}
\centering
\begin{tabular}{p{2.0cm}p{2.5cm}p{2.5cm}}
\toprule
\textbf{Dimension} & \textbf{GraphRAG} & \textbf{RMCA} \\
\midrule
Graph construction & Community detection over document chunks & Typed relational graph from conversation turns \\
Query strategy & Single query against pre-computed summaries & Runtime decomposition into targeted subqueries \\
Token budget & Fixed community summary size & Configurable budget with priority ordering \\
Verbatim evidence & Not available & Verbatim transcript retrieval lane \\
Temporal validity & Not modelled & First-class \texttt{valid\_from}/\texttt{valid\_to} on every node \\
\bottomrule
\end{tabular}
\end{table}


\section{Evaluation}

\subsection{LongMemEval-S}

\textbf{Dataset.} LongMemEval-S~\cite{wu2024longmemeval} is a 500-question benchmark for
multi-session conversational memory retrieval. Each question requires identifying the correct
support session(s) from a haystack of prior sessions. The cleaned \texttt{s} split
(\texttt{xiaowu0162/longmemeval-cleaned}) contains 500 questions with gold support sets of
cardinality 1--6.

\textbf{Metrics.} We report R@5 (recall: does the gold set appear anywhere in the top 5?),
Exact@5 (precision: does the top 5 exactly match the gold set?), Exact@10, and Exact@20.

\textbf{Setup.} Sessions are prepared by converting each haystack session to a text block,
normalising timestamps to UTC ISO format, and embedding unique session texts once with
\texttt{all-MiniLM-L6-v2}. A prepared-session cache keyed by dataset SHA-256, mode, limit,
and embedding model is written on the first run and reused on subsequent runs. All reported
numbers use the warm cache.

\textbf{Modes.} Two retrieval modes are evaluated:
\begin{itemize}
\item \texttt{graph\_raw}: user turns only; weighted score $= 0.72 \times \text{semantic}
+ 0.18 \times \text{lexical} + 0.10 \times \text{temporal}$; top-5 returned directly.
\item \texttt{graph\_hybrid}: user and assistant turns; same raw ranking followed by heuristic
reranking of top-20 using RRF fusion, lexical overlap, and temporal recency bias.
\end{itemize}

\textbf{Results.} Table~\ref{tab:longmemeval} shows the main results.
\texttt{graph\_raw} achieves 89.0\% Exact@5 with no tunable reranking heuristics, making it
the most reproducible baseline. \texttt{graph\_hybrid} scores lower on Exact@5 but recovers
strongly at higher cutoffs (Exact@20 = 98.0\%), indicating it finds all correct sessions but
does not always pack them into the top 5 slots for high-cardinality queries.

\begin{table}[htbp]
\caption{LongMemEval-S Results (500 questions, \texttt{all-MiniLM-L6-v2})}
\label{tab:longmemeval}
\centering
\begin{tabular}{lrrrr}
\toprule
\textbf{Mode} & \textbf{R@5} & \textbf{Exact@5} & \textbf{Exact@10} & \textbf{Exact@20} \\
\midrule
\texttt{graph\_raw}    & \textbf{97.4\%} & \textbf{89.0\%} & 89.0\% & 89.0\% \\
\texttt{graph\_hybrid} & 97.0\%          & 87.2\%          & 94.8\% & \textbf{98.0\%} \\
\bottomrule
\end{tabular}
\end{table}

\textbf{Cardinality breakdown.} Table~\ref{tab:cardinality} shows \texttt{graph\_raw}
performance by gold-set cardinality. R@5 = 100\% at cardinality $\geq 3$ confirms the
embedding reliably finds the correct base session. The Exact@5 drop at high cardinality is
a packing problem: fitting 4--6 chunks into 5 slots while excluding noise is structurally
harder than retrieval itself.

\begin{table}[htbp]
\caption{\texttt{graph\_raw} Performance by Gold-Set Cardinality}
\label{tab:cardinality}
\centering
\begin{tabular}{rrrr}
\toprule
\textbf{Cardinality} & \textbf{$n$} & \textbf{R@5} & \textbf{Exact@5} \\
\midrule
1 & 176 & 95.5\% & 95.5\% \\
2 & 250 & 98.0\% & 93.2\% \\
3 &  41 & 100.0\% & 73.2\% \\
4 &  19 & 100.0\% & 47.4\% \\
5 &  11 & 100.0\% & 45.5\% \\
6 &   3 & 100.0\% &  0.0\% \\
\bottomrule
\end{tabular}
\end{table}

\textbf{Overfitting check.} To verify the scores are not artefacts of the specific
500-question distribution, we ran a held-out split (50 dev / 450 test, seed 42).
Table~\ref{tab:heldout} shows the results. \texttt{graph\_raw} generalises cleanly.
The 5.3\,pp gap for \texttt{graph\_hybrid} is expected: the heuristic reranking weights
(RRF fusion, lexical/coverage scoring) are partially tuned to this distribution.
\texttt{graph\_hybrid} Exact@5 should be treated as an exploratory number rather than a
headline result.

\begin{table}[htbp]
\caption{Held-Out Split Overfitting Check (seed 42)}
\label{tab:heldout}
\centering
\begin{tabular}{lrrr}
\toprule
\textbf{Mode} & \textbf{Dev Exact@5} & \textbf{Test Exact@5} & \textbf{Gap} \\
              & \textbf{($n$=50)}    & \textbf{($n$=450)}    & \\
\midrule
\texttt{graph\_raw}    & 88.0\% & 89.1\% & $+1.1$\,pp --- stable \\
\texttt{graph\_hybrid} & 92.0\% & 86.7\% & $-5.3$\,pp --- dev-set sensitivity \\
\bottomrule
\end{tabular}
\end{table}

\subsection{RLM-Style Synthetic Benchmark}

To evaluate RMCA's behaviour across task types that vary in structural complexity, we
constructed a synthetic benchmark suite following the task taxonomy of the RLM
paper~\cite{zhang2026rlm}. Each family maps a canonical long-context task to a Waggle
memory graph scenario.

\textbf{Task families.} Table~\ref{tab:families} describes the five task families.

\begin{table}[htbp]
\caption{Synthetic Benchmark Task Families}
\label{tab:families}
\centering
\begin{tabular}{p{2.0cm}p{1.8cm}p{3.2cm}}
\toprule
\textbf{Family} & \textbf{RLM equiv.} & \textbf{What it tests} \\
\midrule
S-NIAH-style & RULER S-NIAH~\cite{hsieh2024ruler} & Retrieve one fact from $N$ distractors --- O(1) \\
BrowseComp-Plus-style & BrowseComp-Plus & Multi-hop QA across 3 linked nodes \\
OOLONG-style & OOLONG~\cite{bertsch2025oolong} & Aggregate $N$ entries --- O($n$) \\
OOLONG-Pairs-style & OOLONG-Pairs & Pairwise conflict reasoning --- O($n^2$) \\
CodeQA-style & LongBench-v2 CodeQA~\cite{bai2025longbench} & Codebase/repo understanding \\
\bottomrule
\end{tabular}
\end{table}

\textbf{Methods compared.}
\begin{itemize}
\item \texttt{raw\_context}: dump all graph nodes into the context window (no retrieval).
\item \texttt{query\_graph}: single semantic query, top-$k$ results returned.
\item \texttt{build\_context}: full RMCA pipeline.
\end{itemize}

\textbf{Scales.} Each family is evaluated at $N = 128$, 512, and 2048 nodes.

\textbf{Results.} Table~\ref{tab:rlm_main} shows the main results.

\begin{table}[htbp]
\caption{RLM-Style Synthetic Benchmark Results}
\label{tab:rlm_main}
\centering
\begin{tabular}{lrrrr}
\toprule
\textbf{Family} & \textbf{Scale} & \textbf{raw} & \textbf{query} & \textbf{build} \\
\midrule
S-NIAH       & 128  & 1.000 & 1.000 & \textbf{1.000} \\
S-NIAH       & 512  & 1.000 & 1.000 & \textbf{1.000} \\
S-NIAH       & 2048 & 1.000 & 1.000 & \textbf{1.000} \\
\midrule
BrowseComp+  & 128  & 1.000 & 1.000 & \textbf{1.000} \\
BrowseComp+  & 512  & 1.000 & 1.000 & \textbf{1.000} \\
BrowseComp+  & 2048 & 1.000 & 1.000 & \textbf{1.000} \\
\midrule
OOLONG-Pairs & 128  & 0.000 & 0.000 & \textbf{1.000} \\
OOLONG-Pairs & 512  & 0.000 & 0.000 & \textbf{1.000} \\
OOLONG-Pairs & 2048 & 0.000 & 0.000 & \textbf{1.000} \\
\midrule
CodeQA       & 128  & 0.000 & 1.000 & \textbf{1.000} \\
CodeQA       & 512  & 0.000 & 1.000 & \textbf{1.000} \\
CodeQA       & 2048 & 0.000 & 1.000 & \textbf{1.000} \\
\midrule
OOLONG       & 128  & 0.885 & 0.242 & 0.513 \\
OOLONG       & 512  & 0.403 & 0.035 & 0.224 \\
OOLONG       & 2048 & 0.000 & 0.010 & 0.069 \\
\bottomrule
\end{tabular}
\end{table}

\textbf{Token efficiency.} On S-NIAH tasks, \texttt{build\_context} uses 13--14\% of the
tokens consumed by \texttt{raw\_context} while achieving the same score (7--8$\times$
reduction). On OOLONG-Pairs tasks, \texttt{build\_context} uses $\approx$36\% of
\texttt{raw\_context} tokens while achieving 1.0 vs.\ 0.0 score.

\textbf{Key findings.}
\begin{enumerate}
\item \textbf{Pairwise conflict reasoning (OOLONG-Pairs).} Both \texttt{raw\_context} and
\texttt{query\_graph} score 0.0 at all scales. \texttt{build\_context} scores 1.0 at all
scales. RMCA succeeds on this task because decomposition generates targeted conflict/constraint subqueries and the formatted context pack exposes the conflict relation in an answerable structure. Current ablations isolate decomposition---not graph expansion or conflict resolution---as the load-bearing component. Typed edges and conflict resolution are part of the RMCA design, but current synthetic benchmarks do not yet isolate their causal contribution.
\item \textbf{Codebase understanding (CodeQA).} \texttt{raw\_context} scores 0.0 at all
scales (the relevant code fact is buried in a large context dump). \texttt{query\_graph}
and \texttt{build\_context} both score 1.0, confirming that targeted retrieval is sufficient
for this task type.
\item \textbf{Simple factual retrieval (S-NIAH, BrowseComp-Plus).} All three methods score
1.0 at all scales. RMCA uses 7--8$\times$ fewer tokens than raw context while achieving the
same score.
\item \textbf{Linear aggregation (OOLONG).} \texttt{raw\_context} scores highest at small
scale (0.885 at $N$=128) but degrades to 0.0 at $N$=2048 as the context window fills.
\texttt{build\_context} scores 0.513 at $N$=128 and degrades to 0.069 at $N$=2048. This
exposes a fundamental limitation: tasks requiring aggregation over all $N$ entries cannot
be fully satisfied by any fixed-budget retrieval system.
\end{enumerate}

\subsection{Latency}

Local latency on SQLite with deterministic embeddings (warm cache):

\begin{table}[htbp]
\caption{Operation Latency (SQLite, warm cache)}
\label{tab:latency}
\centering
\begin{tabular}{lr}
\toprule
\textbf{Operation} & \textbf{Mean latency} \\
\midrule
\texttt{observe\_conversation} & 1.54\,ms \\
\texttt{query\_graph}          & 1.60\,ms \\
\texttt{graph\_diff}           & 0.80\,ms \\
\texttt{build\_context} ($N$=128)  & $\approx$22\,ms \\
\texttt{build\_context} ($N$=2048) & $\approx$258--415\,ms \\
\bottomrule
\end{tabular}
\end{table}

\texttt{build\_context} latency scales with graph size due to the multi-subquery retrieval
loop. At $N$=2048, latency is in the 250--415\,ms range, which is acceptable for agent use
but not for real-time interactive applications.


\subsection{LLM-Backed Answer-Level Evaluation}

To validate that retrieval-level scores translate to answer quality, we conducted an
LLM-backed answer-level evaluation using Groq \texttt{llama-3.3-70b-versatile} as the
answering model. Each method's context pack was injected into a structured prompt; the
model's answer was evaluated for exact match (EM), F1, and hallucination rate.

\textbf{Setup.} Evaluation was run across 3 seeds (42, 43, 44) and 2 scales (128, 512 nodes)
on the pairwise conflict task. Results were consistent across all seeds and scales.

\textbf{Results.} Table~\ref{tab:groq_pairwise} shows the pairwise task results.

\begin{table}[htbp]
\caption{LLM Answer-Level Evaluation --- Pairwise Conflict Task\\
(Groq \texttt{llama-3.3-70b-versatile}, seeds 42/43/44, scales 128/512)}
\label{tab:groq_pairwise}
\centering
\begin{tabular}{lrrrr}
\toprule
\textbf{Method} & \textbf{EM} & \textbf{F1} & \textbf{Hall.} & \textbf{Tokens} \\
\midrule
\texttt{rmca\_full}  & \textbf{1.00} & \textbf{1.00} & \textbf{0.00} & 435--515 \\
\texttt{query\_graph} & 0.00 & 0.00 & 0.00 & 98 \\
\texttt{bm25\_topk}  & 0.00 & 0.00--0.12 & 0.00--1.00 & 1390--1422 \\
\texttt{raw\_context} & 0.00 & 0.00--0.12 & 0.00--1.00 & 1390--1422 \\
\texttt{hybrid\_rrf} & 0.00 & 0.00--0.12 & 0.00--1.00 & 1389--1422 \\
\bottomrule
\end{tabular}
\end{table}

\textbf{Key findings.}
\begin{enumerate}
\item RMCA is the only method achieving EM=1.0, F1=1.0, and hallucination rate=0.0 on
pairwise conflict tasks, consistent across all 3 seeds and both scales.
\item \texttt{query\_graph} retrieves insufficient context (98 tokens) and the model
correctly marks the question as \texttt{insufficient\_context} rather than hallucinating.
\item Flat baselines (\texttt{bm25\_topk}, \texttt{raw\_context}, \texttt{hybrid\_rrf})
inject 1390--1422 tokens but either fail to answer correctly or hallucinate, because the
conflict annotation is not surfaced in the flat context dump.
\item The pattern is consistent across seeds 42, 43, 44 and scales 128, 512, confirming
the result is not a random artefact.
\end{enumerate}

\textbf{Caveat.} These results use one answering model and should be replicated across
models. The model is an answerer, not an independent human judge. The
\texttt{DeterministicAnswerer} is a reproducible lower-bound evaluator; its scores are not
equivalent to human preference ratings or LLM-judge quality assessments.

\subsection{Ablation Study}

To identify which RMCA components are load-bearing, we ran a systematic ablation study
across 7 variants on the \texttt{pairwise\_hidden\_edge} benchmark (conflict represented
only by typed \texttt{contradicts} edges; no conflict vocabulary in node content).

\textbf{Setup.} 4 variants $\times$ 2 scales (128, 512) $\times$ 3 seeds (42, 43, 44).
All seeds produced identical scores (deterministic benchmark).

\textbf{Results.} Table~\ref{tab:ablation} shows the ablation results.

\begin{table}[htbp]
\caption{RMCA Ablation Study --- \texttt{pairwise\_hidden\_edge} Benchmark\\
(scales 128/512, seeds 42/43/44; all seeds identical)}
\label{tab:ablation}
\centering
\begin{tabular}{lrrrp{2.0cm}}
\toprule
\textbf{Variant} & \textbf{S=128} & \textbf{S=512} & \textbf{$\Delta$} & \textbf{Status} \\
\midrule
\texttt{rmca\_full}               & 1.000 & 1.000 & ---    & Baseline \\
\texttt{rmca\_no\_decomposition}  & 0.000 & 0.000 & $-1.000$ & \textbf{Load-bearing} \\
\texttt{rmca\_no\_graph\_expansion} & 1.000 & 1.000 & 0.000 & Not isolated \\
\texttt{rmca\_no\_conflict\_resolution} & 1.000 & 1.000 & 0.000 & Not isolated \\
\bottomrule
\end{tabular}
\end{table}

\textbf{Real-embedding confirmation.} The ablation was repeated with Waggle's production
\texttt{all-MiniLM-L6-v2} model (384-dim sentence-transformers) to rule out artefacts of
the deterministic hash-based embedding. Table~\ref{tab:ablation_real} shows the results.

\begin{table}[htbp]
\caption{Ablation with Real Embeddings (\texttt{all-MiniLM-L6-v2})\\
vs.\ Deterministic Embeddings}
\label{tab:ablation_real}
\centering
\begin{tabular}{lrrrr}
\toprule
\textbf{Variant} & \textbf{Det.} & \textbf{Real} & \textbf{Diff.} \\
\midrule
\texttt{rmca\_full}               & 1.000 & 1.000 & 0.000 \\
\texttt{rmca\_no\_graph\_expansion} & 1.000 & 1.000 & 0.000 \\
\texttt{rmca\_no\_conflict\_resolution} & 1.000 & 1.000 & 0.000 \\
\texttt{rmca\_no\_decomposition}  & 0.000 & 0.000 & 0.000 \\
\bottomrule
\end{tabular}
\end{table}

\textbf{Root cause analysis.} The real embedding model produces identical results to the
deterministic model because the root cause is \emph{benchmark construction}, not the
embedding model. The \texttt{pairwise\_hidden\_edge} node labels (e.g., ``Cloud database'',
``Local deployment required'') are semantically distinctive enough that
\texttt{all-MiniLM-L6-v2} retrieves the relevant choice and constraint nodes directly via
cosine similarity, without needing to traverse the \texttt{contradicts} edge.

To causally isolate graph expansion and conflict resolution, the benchmark would need nodes
where: (1) choice and constraint labels are semantically similar to distractors, so that
semantic retrieval alone cannot distinguish them; and (2) the only signal distinguishing the
gold conflict pair is the \texttt{contradicts} edge. This requires a fundamentally different
benchmark design.

\textbf{Summary.} Decomposition is the only RMCA component causally isolated as
load-bearing on the current synthetic benchmarks. Graph expansion and conflict resolution
may be load-bearing in real-world settings, but the current benchmark cannot demonstrate
this due to construction limitations.

\subsection{ContextReset Benchmark}

The ContextReset benchmark evaluates whether a memory system can correctly reconstruct
project state after a context window reset. It tests six fields: decision recall,
constraint recall, next-step accuracy, superseded context handling, active decision
preference, and evidence coverage.

\textbf{Setup.} Scale $N$=128, after the decomposition fix (short continuation queries
now use broad project-state subqueries).

\textbf{Results.} Table~\ref{tab:context_reset} shows the results.

\begin{table}[htbp]
\caption{ContextReset Benchmark Results ($N$=128)}
\label{tab:context_reset}
\centering
\begin{tabular}{lrr}
\toprule
\textbf{Method} & \textbf{Score} & \textbf{Tokens} \\
\midrule
\texttt{rmca\_full}  & \textbf{1.000} & 315 \\
\texttt{query\_graph} & 0.875 & 96 \\
\texttt{raw\_context} & 0.000 & 1413 \\
\texttt{bm25\_topk}  & 0.000 & 1413 \\
\texttt{no\_memory}  & 0.000 & 0 \\
\bottomrule
\end{tabular}
\end{table}

\texttt{rmca\_full} achieves a perfect score of 1.000 across all six evaluation fields.
\texttt{query\_graph} scores 0.875, missing one field due to insufficient context breadth.
Both \texttt{raw\_context} and \texttt{bm25\_topk} score 0.0 despite injecting 1413 tokens,
because the flat context dump does not surface the structured project state needed to answer
the multi-field question. \texttt{no\_memory} scores 0.0 as expected.

The key insight from this benchmark is that structured context assembly (RMCA) outperforms
both flat retrieval and raw context dump on multi-aspect project-state reconstruction tasks,
using only 22\% of the tokens consumed by the flat baselines.

\subsection{Comparison with RLM Paper}

The RLM paper~\cite{zhang2026rlm} reports results on the real OOLONG and OOLONG-Pairs
datasets using frontier models. Table~\ref{tab:rlm_compare} shows the comparison.

\textbf{Note: Rows in this table are not a leaderboard.} The RLM paper evaluates real long-context datasets (OOLONG \texttt{trec\_coarse}, OOLONG-Pairs) with frontier models (GPT-5, Qwen3-Coder-480B). Waggle evaluates synthetic graph-memory analogues with a 70B open model. Numbers are not comparable. The shared qualitative pattern is that structured decomposition/assembly outperforms flat/base methods on pairwise/conflict-style tasks.

\begin{table}[htbp]
\caption{Qualitative Alignment with RLM-Style Findings\\
(Not Numerically Comparable --- Not a Leaderboard)}
\label{tab:rlm_compare}
\centering
\begin{tabular}{p{2.8cm}p{1.5cm}p{1.5cm}p{1.2cm}}
\toprule
\textbf{System} & \textbf{OOLONG} & \textbf{OOLONG-Pairs F1} & \textbf{Dataset} \\
\midrule
RLM (GPT-5)~\cite{zhang2026rlm} & $\approx$68\% & 58.0\% & Real \texttt{trec\_coarse} \\
RLM (Qwen3-Coder-480B)~\cite{zhang2026rlm} & $\approx$73\% & 23.1\% & Real \texttt{trec\_coarse} \\
Base models~\cite{zhang2026rlm} & --- & $<$0.1\% & Real \texttt{trec\_coarse} \\
\midrule
Waggle RMCA (scoped) & 55\% & --- & Synthetic (20 cases) \\
Waggle RMCA (scoped fullscan) & 100\% & --- & Synthetic (20 cases) \\
Waggle RMCA (pairwise) & --- & \textbf{100\%} & Synthetic (9 nodes) \\
\bottomrule
\end{tabular}
\end{table}

\textbf{Important caveats.}
\begin{enumerate}
\item The Waggle synthetic OOLONG results (55\% scoped, 100\% scoped-fullscan) are
\emph{not comparable} to the RLM paper's 68--73\% on the real \texttt{trec\_coarse} split.
The datasets, models, and scales differ fundamentally.
\item The OOLONG real-30 run was interrupted by API token-per-day limits before completion;
the 20-case synthetic result is preliminary.
\item The RMCA pairwise F1 of 100\% is on a synthetic task with 9 nodes and explicit
\texttt{contradicts} edges, which is substantially simpler than the real OOLONG-Pairs split.
\item The qualitative pattern is consistent: structured assembly outperforms flat retrieval
on conflict-aware tasks. Base models score $<$0.1\% F1 on OOLONG-Pairs while structured
recursive assembly recovers to 23--58\% F1 in the RLM paper; RMCA achieves 100\% F1 on
the synthetic pairwise task.
\end{enumerate}


\subsection{Supported and Unsupported Claims}
\label{sec:claims}

Table~\ref{tab:claims} summarises which claims are supported by current evidence
and which require further work. This table is intended to help reviewers assess
the scope of our contributions.

\begin{table*}[htbp]
\caption{Supported and Unsupported Claims}
\label{tab:claims}
\centering
\begin{tabular}{p{3.5cm}p{1.8cm}p{4.5cm}p{3.5cm}}
\toprule
\textbf{Claim} & \textbf{Status} & \textbf{Evidence} & \textbf{Caveat} \\
\midrule
RMCA improves pairwise answerability & \textbf{Supported} & Groq eval: EM/F1=1.0, hall=0.0 across 3 seeds $\times$ 2 scales & Synthetic pairwise task; one answer model \\
\midrule
Decomposition is load-bearing & \textbf{Supported} & \texttt{no\_decomposition} and \texttt{random\_subqueries} drop 1.0$\to$0.0 & Current synthetic tasks only \\
\midrule
RMCA helps ContextReset & \textbf{Supported} & \texttt{rmca\_full}=1.0 vs \texttt{query\_graph}=0.875 and raw/bm25/no\_memory=0.0 & Synthetic project-state benchmark \\
\midrule
RMCA reduces tokens vs raw context & \textbf{Supported (selected tasks)} & S-NIAH: 13--14\% of raw tokens; pairwise: $\approx$31--38\% & Not universally lower than \texttt{query\_graph} \\
\midrule
Graph expansion is load-bearing & \textbf{Not supported yet} & Ablation shows no delta at scales 128/512 & Benchmark construction may not isolate edge traversal \\
\midrule
Conflict resolution is load-bearing & \textbf{Not supported yet} & Ablation shows no delta at scales 128/512 & Direct retrieval already surfaces conflict nodes \\
\midrule
RMCA solves linear aggregation & \textbf{Not supported} & OOLONG-style degrades with scale & O($n$) coverage under fixed budget \\
\midrule
Results numerically comparable to RLM paper & \textbf{Not supported} & Different datasets, models, scales & Qualitative comparison only \\
\midrule
Results generalise to real-world traces & \textbf{Not supported yet} & No real trace benchmark & Future work \\
\bottomrule
\end{tabular}
\end{table*}


\section{Discussion}

\subsection{When RMCA Helps}

RMCA provides the largest benefit over raw context and single-query retrieval on tasks
that require:
\begin{itemize}
\item \textbf{Conflict and contradiction detection.} On current pairwise benchmarks, query decomposition is the primary causally isolated load-bearing component. Disabling decomposition or replacing it with random subqueries drops score from 1.0 to 0.0. Disabling graph expansion or explicit conflict resolution does not reduce score, because direct retrieval already surfaces the relevant conflict nodes in this synthetic setup. Typed edges and conflict resolution are part of the RMCA design, but current synthetic benchmarks do not yet isolate their causal contribution.
\item \textbf{Multi-aspect information needs.} Query decomposition (Step 1) allows RMCA
to retrieve memory relevant to decisions, constraints, implementation details, and conflicts
in a single call, where a single-query retrieval would only surface the most semantically
similar nodes.
\item \textbf{Token-constrained contexts.} At large graph sizes, raw context dumps exceed
any practical context window. RMCA's token-budget compression (Step 7) ensures the context
pack fits regardless of graph size.
\item \textbf{Project-state reconstruction.} The ContextReset benchmark demonstrates that
RMCA's structured assembly correctly reconstructs multi-field project state (decisions,
constraints, next steps, superseded context) where flat retrieval fails.
\end{itemize}

\subsection{When RMCA Does Not Help}

RMCA does not improve over simpler methods on:
\begin{itemize}
\item \textbf{Simple factual retrieval.} When the answer is a single node with high
semantic similarity to the query, \texttt{query\_graph} is sufficient and faster.
\item \textbf{Linear aggregation.} Tasks requiring enumeration of all $N$ entries are
fundamentally O($n$) information needs. No fixed-budget retrieval system can guarantee
full coverage. RMCA degrades more gracefully than raw context at large scale but does not
solve this problem.
\end{itemize}

\subsection{Limitations}

\begin{enumerate}
\item \textbf{Synthetic benchmark data.}
The RLM-style benchmark uses deterministic synthetic Waggle memory tasks. Numerical results
must not be compared to the RLM paper or other long-context benchmarks until the exact
public datasets are run with a matching model setup. The OOLONG real-30 run was interrupted
by API token-per-day limits before completion; the 20-case synthetic result (55\% scoped,
100\% scoped-fullscan) is not comparable to the RLM paper's 68--73\% on the real
\texttt{trec\_coarse} split.

\item \textbf{Token estimation is approximate.}
$\widehat{\text{tok}}(s) = \lfloor |s|/4 \rfloor$ may differ from actual tokenizer counts
by $\pm$20\%.

\item \textbf{Decomposition is heuristic, not learned.} Subquery generation uses keyword
pattern matching and does not use an LLM. It may produce suboptimal decompositions for
queries outside the project/coding and generic-memory categories.

\item \textbf{Graph expansion is bounded.} Expansion is limited to 2 hops from the top-3
seed nodes. Deep reasoning chains requiring more than 2 hops are not guaranteed to be
surfaced.

\item \textbf{Graph expansion and conflict resolution are not yet causally isolated.}
Ablation studies on both the standard \texttt{pairwise} and the harder
\texttt{pairwise\_hidden\_edge} benchmark (conflict represented only by typed
\texttt{contradicts} edges, no conflict vocabulary in node content) show zero delta when
graph expansion or conflict resolution is disabled, using both the deterministic hash-based
embedding model and Waggle's production \texttt{all-MiniLM-L6-v2} model. The root cause is
benchmark construction: node labels are semantically distinctive enough that direct
retrieval surfaces conflict nodes without edge traversal. Isolating these components
requires a benchmark with semantically indistinguishable node labels and a scorer that
requires explicit conflict annotation in the context pack.

\item \textbf{LLM answer-level evaluation uses a single model.}
The Groq \texttt{llama-3.3-70b-versatile} answer-level results use one answering model and
should be replicated across models. A second-model run with \texttt{llama-3.1-8b-instant}
was deferred because \texttt{GROQ\_API\_KEY} was not available in the evaluation environment.
The model is an answerer, not an independent human judge. The
\texttt{DeterministicAnswerer} is a reproducible lower-bound evaluator; its scores are not
equivalent to human preference ratings or LLM-judge quality assessments.

\item \textbf{LongMemEval scope.} The adapter evaluates session retrieval only. End-to-end
graph extraction quality, contradiction handling across sessions, and context-bundle
usefulness for model handoff are not yet measured.

\item \textbf{\texttt{graph\_hybrid} reranking sensitivity.} The 5.3\,pp dev/test gap
confirms the hybrid reranking weights are partially tuned to this distribution and should
not be treated as a portable headline number.

\item \textbf{OOLONG-Pairs comparison is qualitative, not quantitative.}
The RLM paper reports RLM(GPT-5) at 58.0\% F1 and RLM(Qwen3-Coder-480B) at 23.1\% F1 on
the real OOLONG-Pairs split, while base models score $<$0.1\%. RMCA achieves 100\% F1 on
the synthetic pairwise task (9 nodes, explicit conflict edges). The qualitative pattern is
consistent---structured assembly outperforms flat retrieval on conflict-aware tasks---but
the synthetic task is substantially simpler than the real OOLONG-Pairs split and the
underlying models differ.
\end{enumerate}

\section{Conclusion}

We presented Waggle, a local-first persistent memory system for AI agents, and RMCA,
an algorithm for assembling compact, high-signal context packs from a typed knowledge graph.

On LongMemEval-S, \texttt{graph\_raw} achieves 89.0\% Exact@5 with no overfitting
confirmed by a held-out split (+1.1\,pp gap), demonstrating strong session retrieval
on a real multi-session benchmark.

On synthetic pairwise conflict reasoning and ContextReset tasks, RMCA achieves perfect
scores while flat baselines score 0.0. LLM-backed answer-level evaluation confirms
EM=1.0, F1=1.0, and hallucination rate=0.0 for RMCA on pairwise tasks across 3 seeds
and 2 scales. Ablation studies identify \textbf{query decomposition as the only currently
causally isolated load-bearing component}: disabling it drops pairwise score from 1.0 to 0.0.
Graph expansion and conflict resolution remain plausible contributors to RMCA's design,
but current synthetic benchmarks do not yet isolate their causal contribution due to
benchmark construction limitations.

On linear aggregation tasks, RMCA degrades gracefully but exposes a fundamental O($n$)
coverage limit shared by all fixed-budget retrieval systems.

The key insight is that agent memory is not a retrieval problem alone---it is a
retrieval-plus-reasoning problem. Retrieving the right facts is necessary but not sufficient;
the memory system must also surface which facts are current, which are superseded, and which
are in tension with each other. RMCA provides the architectural machinery for this, but
causal isolation of individual components remains an open empirical question.

Future work includes: (1) replacing heuristic query decomposition with a learned decomposer;
(2) evaluating on the full RLM public benchmark suite with a matching model setup;
(3) measuring end-to-end graph extraction quality on LongMemEval;
(4) hardening the hybrid reranking weights against distribution shift;
(5) constructing a \texttt{pairwise\_hidden\_edge} variant with semantically indistinguishable
node labels to causally isolate graph expansion and conflict resolution; and
(6) evaluating on real agent traces, which is the most important next step for
establishing external validity.


\begin{thebibliography}{99}

\bibitem{lewis2020rag}
P.~Lewis, E.~Perez, A.~Piktus, F.~Petroni, V.~Karpukhin, N.~Goyal, H.~K\"{u}ttler,
M.~Lewis, W.-t. Yih, T.~Rockt\"{a}schel, S.~Riedel, and D.~Kiela,
``Retrieval-augmented generation for knowledge-intensive {NLP} tasks,''
in \emph{Advances in Neural Information Processing Systems (NeurIPS)}, 2020.

\bibitem{edge2024graphrag}
D.~Edge, H.~Trinh, N.~Cheng, J.~Bradley, A.~Chao, A.~Mody, S.~Truitt, and J.~Larson,
``From local to global: A {Graph RAG} approach to query-focused summarization,''
\emph{arXiv preprint arXiv:2404.16130}, 2024.

\bibitem{packer2023memgpt}
C.~Packer, S.~Wooders, K.~Lin, V.~Fang, S.~I.~Patil, I.~Stoica, and J.~E.~Gonzalez,
``{MemGPT}: Towards {LLMs} as operating systems,''
\emph{arXiv preprint arXiv:2310.08560}, 2023.

\bibitem{park2023generative}
J.~S.~Park, J.~C.~O'Brien, C.~J.~Cai, M.~R.~Morris, P.~Liang, and M.~S.~Bernstein,
``Generative agents: Interactive simulacra of human behavior,''
in \emph{Proc.\ ACM Symposium on User Interface Software and Technology (UIST)}, 2023.

\bibitem{wu2024longmemeval}
X.~Wu, H.~Wang, W.~Xiong, C.~Hu, and W.~Y.~Wang,
``{LongMemEval}: Benchmarking chat assistants on long-term interactive memory,''
\emph{arXiv preprint arXiv:2410.10813}, 2024.

\bibitem{reimers2019sbert}
N.~Reimers and I.~Gurevych,
``Sentence-{BERT}: Sentence embeddings using {Siamese BERT}-networks,''
in \emph{Proc.\ Conference on Empirical Methods in Natural Language Processing (EMNLP)},
2019.

\bibitem{zhang2026rlm}
A.~L.~Zhang, T.~Kraska, and O.~Khattab,
``Recursive language models,''
\emph{arXiv preprint arXiv:2512.24601}, 2026.

\bibitem{bertsch2025oolong}
A.~Bertsch, A.~Pratapa, T.~Mitamura, G.~Neubig, and M.~R.~Gormley,
``{OOLONG}: Evaluating long context reasoning and aggregation capabilities,''
\emph{arXiv preprint arXiv:2511.02817}, 2025.

\bibitem{hsieh2024ruler}
C.-P.~Hsieh, S.~Sun, S.~Kriman, S.~Acharya, D.~Rekesh, F.~Jia, and B.~Ginsburg,
``{RULER}: What's the real context size of your long-context language models?''
\emph{arXiv preprint arXiv:2404.06654}, 2024.

\bibitem{bai2025longbench}
Y.~Bai, X.~Lv, J.~Zhang, Y.~He, J.~Qi, L.~Hou, J.~Tang, Y.~Dong, and J.~Li,
``{LongBench v2}: Towards deeper understanding and reasoning on realistic long-context
multitasks,''
\emph{arXiv preprint arXiv:2412.15204}, 2025.

\end{thebibliography}

\appendix

\section{RMCA Pseudocode}

\begin{enumerate}
\item \textbf{DECOMPOSE:} Classify query intent; generate $\leq k$ subqueries $\{sq_i\}$ with priority weights and retrieval modes.
\item \textbf{RETRIEVE:} $H \leftarrow \bigcup_i \text{retrieve}(sq_i, G, T)$ --- issue each subquery against graph and transcript store.
\item \textbf{EXPAND:} seeds $\leftarrow$ top-5 nodes in $H$; for each $s \in$ seeds[:3]: $H \leftarrow H \cup \text{graph\_hop}(s, G, \min(d,2))$.
\item \textbf{RESOLVE:} For \texttt{updates} edges: penalise target ($\times 0.3$), boost source ($+0.15$). For \texttt{contradicts} edges: record conflict pair. For expired nodes: penalise ($\times 0.2$).
\item \textbf{DEDUPLICATE:} $H \leftarrow \{\arg\max_{\text{score}} h : h.\text{node\_id}\}$ --- keep highest-scored copy per node.
\item \textbf{RANK:} Sort $H$ by (is\_superseded $\uparrow$, score $\downarrow$, label $\uparrow$).
\item \textbf{COMPRESS:} Pack sections [decisions, constraints, impl., unfinished, conflicts, superseded, evidence] until $\widehat{\text{tok}}(C) > B \times 1.15$.
\item \textbf{FORMAT:} Prepend structured header; attach provenance and conflict annotations.
\item[] \textbf{Return} $C$, provenance($H$), conflict\_entries.
\end{enumerate}

\section{Reproduction Commands}

\subsection{LongMemEval-S}

\begin{verbatim}
# Full 500-question run -- graph_raw
python scripts/benchmark_longmemeval.py \
  benchmarks/longmemeval/longmemeval_s_cleaned.json \
  --mode graph_raw \
  --output results/graph_raw.json

# Held-out split (overfitting check)
python scripts/benchmark_longmemeval.py \
  benchmarks/longmemeval/longmemeval_s_cleaned.json \
  --mode graph_raw --held-out \
  --output results/graph_raw_heldout.json
\end{verbatim}

\subsection{RLM-Style Synthetic Benchmark}

\begin{verbatim}
python benchmarks/rlm_style_waggle_eval.py \
  --db /tmp/waggle_rlm_eval \
  --scales 128 512 2048 \
  --methods raw_context query_graph build_context \
  --families sniah multihop pairwise codeqa linear_agg \
  --token-budget 1200 --seed 42 \
  --output benchmark_results/
\end{verbatim}

\subsection{Ablation Study}

\begin{verbatim}
python benchmarks/run_ablation.py \
  --benchmark pairwise_hidden_edge \
  --scales 128 512 \
  --seeds 42 43 44 \
  --real-embeddings \
  --output benchmark_results/ablation_results_real_emb.csv
\end{verbatim}

\end{document}
